\chapter{LSP protocol internals}
\label{app:lsp-protocol}
\index{LSP!protocol internals}\index{JSON-RPC}\index{capability negotiation}\index{LSP!wire format}

This appendix supplements Chapter~\ref{ch:eglot-lsp} with details
of the Language Server Protocol's wire format and capability
negotiation.

\section{JSON-RPC wire format}

LSP communication uses JSON-RPC 2.0 over stdin/stdout.  Each message
has an HTTP-style header followed by a JSON body:

\begin{minted}{text}
Content-Length: 84\r\n
\r\n
{"jsonrpc":"2.0","id":1,"method":"initialize","params":{...}}
\end{minted}

Messages fall into three categories:

\begin{description}
  \item[Request] Client → server (or vice versa); has \texttt{id} and \texttt{method}.  Expects a response.
  \item[Response] Matches a request by \texttt{id}; contains either \texttt{result} or \texttt{error}.
  \item[Notification] Has \texttt{method} but no \texttt{id}.  No response expected.
\end{description}

Eglot serialises and deserialises JSON using the built-in
\texttt{jsonrpc.el} library, which dispatches incoming messages to
handler functions registered per method name.

\section{Capability negotiation}

The \texttt{initialize} request is the first message sent by the
client.  It advertises the client's capabilities (what the editor
supports) and receives the server's capabilities (what intelligence
the server can provide):

\begin{minted}{text}
Client → Server:  initialize { capabilities: { ... } }
Server → Client:  { capabilities: { completionProvider: ..., ... } }
Client → Server:  initialized {}   (notification)
\end{minted}

Eglot reads the server's response and enables only the Emacs
features backed by a declared server capability.  For example, if
the server omits \texttt{renameProvider}, \texttt{eglot-rename}
signals an error immediately rather than sending an unsupported
request.

\section{\texttt{read-process-output-max} and throughput}

The default \texttt{read-process-output-max} (4\,KB on most systems)
limits how much data Emacs reads from a subprocess pipe in a single
\texttt{read} syscall.  LSP servers that return large completion lists
(hundreds of items) can saturate this limit, causing multiple read
cycles per response.  Setting the variable to 1--2\,MB lets Emacs
consume an entire LSP response in one call, reducing completion
latency significantly.

\begin{minted}{elisp}
(setq read-process-output-max (* 2 1024 1024))  ; 2 MB
\end{minted}

This variable affects \emph{all} subprocesses, not just LSP servers.
